\documentclass{article}
\usepackage{graphicx}
\usepackage{amsmath}
\usepackage{amssymb}
\usepackage{geometry}
\usepackage{wrapfig}
\bibliographystyle{plain}

\geometry{a4paper,left=2.4cm,right=2.4cm,top=1.9cm,bottom=1.9cm}

\title{Data-Aware and Budget-Constrained\\LoRA Mixture-of-Experts Allocation for Large Language Models}
\author{Zhi-Quan Feng}
\date{}

\begin{document}
\maketitle
\fontsize{11pt}{11pt}\selectfont

\section{Introduction}


Large language models (LLMs) have become general-purpose foundations for natural language understanding and reasoning. Their broad capability, however, does not remove the need for task adaptation. A pretrained model must still be specialized for domains, datasets, user preferences, or newly emerging tasks. Repeating full-parameter fine-tuning for every use case is increasingly impractical because it requires storing, optimizing, and maintaining billions of task-specific parameters.

Parameter-efficient fine-tuning (PEFT) addresses this problem by keeping most pretrained weights frozen and learning only a small set of additional parameters. LoRA\cite{LORA}, in particular, represents a weight update using two compact matrices and has become a widely used baseline because it substantially reduces trainable parameters and optimizer memory. Nevertheless, parameter efficiency alone does not guarantee that the limited adaptation budget is used efficiently. Conventional LoRA normally assigns the same rank and the same structural form to every selected module. This uniform design implicitly assumes that all layers and projections require comparable adaptation capacity.

That assumption is unlikely to hold across a deep language model. Lower and higher layers play different representational roles; attention and feed-forward projections transform different feature spaces; and the modules that matter for one downstream task may be less important for another. Some modules may require only a small, shared correction, whereas others may need to represent several distinct patterns for different inputs. A uniform adapter therefore creates two possible forms of waste: allocating unnecessary capacity to simple modules and restricting modules whose adaptation demands are more diverse.

Mixture-of-Experts (MoE) models provide a complementary form of conditional capacity. Instead of applying all parameters to every token, a router selects a small subset of experts\cite{SPARSEMOE,HYPERROUTER,DEEPSEEKV3}. Combining low-rank adapters with conditional experts is attractive because it may increase task-specific expressiveness without fine-tuning the entire backbone. Yet applying the same MoE structure everywhere simply replaces one uniform design with another. It can introduce redundant experts, unstable routing, additional memory, and an unfair parameter advantage over simpler baselines.

This project studies a more fundamental question: \emph{how should a fixed parameter-efficient adaptation budget be distributed across the heterogeneous modules of an LLM?} The goal is not to present one adapter architecture in isolation, but to understand when ordinary low-rank adaptation is sufficient, when conditional expert capacity is justified, and how these decisions can be made from limited task data before expensive fine-tuning. Solving this problem would make PEFT more accurate, more resource-aware, and easier to deploy across models and tasks.

\section{Related Works and Innovations}

\subsection{Parameter-Efficient Adaptation and Non-uniform Rank}

For a frozen weight matrix $W\in\mathbb{R}^{d_{\mathrm{out}}\times d_{\mathrm{in}}}$, LoRA learns $W'=W+BA$, where the rank $r$ of the compact matrices controls both adaptation capacity and parameter cost\cite{LORA}. Its simplicity has made LoRA a standard mechanism for maintaining many specialized versions of one foundation model. However, assigning the same rank to every target module assumes that all layers and projections require equal capacity. AdaLoRA challenged this assumption by dynamically allocating a rank budget according to parameter importance\cite{ADALORA}. More recent research continues to investigate dynamic rank use, parameter sharing, and the relation between low-rank optimization and full fine-tuning.

This progression shows that rank is not merely an implementation setting; it determines which subspace is available for learning. Our team's recent work, \emph{Learning in the Fisher Subspace: A Guided Initialization for LoRA Fine-Tuning}, further demonstrates that task-aware sensitivity is more informative than pretrained weight geometry alone when selecting the initial adaptation subspace\cite{FISHERSUBSPACE}. This result motivates the present proposal to move beyond choosing directions within a fixed adapter and study how task-conditioned capacity should be distributed across the model itself.

\subsection{Mixtures of Low-Rank Experts}

Sparse MoE architectures increase conditional capacity by routing each input to only a subset of experts\cite{SPARSEMOE,HYPERROUTER,DEEPSEEKV3}. Recent PEFT work has brought this principle into low-rank adaptation. AdaMoLE adapts expert activation to input complexity rather than relying on a fixed top-$k$ rule\cite{ADAMOLE}. GOAT, published at ICML 2025, studies SVD-structured LoRA experts and the optimization mismatch between LoRA--MoE and full fine-tuning\cite{GOAT}. D-MoLE, also published at ICML 2025, shows that layer-wise expert allocation under a controlled budget is important when task demands evolve\cite{DMOLE}. HILO jointly considers expert number and adapter rank across layers\cite{HILO}, while sensitivity-driven LoRA-SMoE uses limited task data to allocate experts under a resource constraint\cite{LORASMOE}. DR-LoRA further reports that uniform expert ranks can under-provision task-relevant experts while wasting parameters on less relevant ones\cite{DRLORA}.

Taken together, these studies confirm that adaptive low-rank experts are an important and rapidly developing direction. They also expose an unresolved issue: expert selection, rank allocation, routing sparsity, and parameter accounting are often optimized separately or for a specialized setting. A method can appear efficient while storing more parameters, activating more rank, or relying on expensive gradients during architecture selection. There remains no settled principle for deciding, before full training, which modules should remain ordinary LoRA modules, which genuinely benefit from experts, and how a common model-wide budget should be divided among modules of unequal size.

\subsection{Why This Topic Must Be Investigated}

The importance of this problem follows from the way LLMs are now used. Organizations increasingly maintain one pretrained backbone and many domain- or task-specific adapters. As the number of tasks grows, even ``parameter-efficient'' adapters become a substantial storage, training, and serving burden. Recent systems research has consequently begun to study how thousands of LoRA adapters can be compressed and served with little overhead\cite{COMPRESSLORA}. Better allocation at training time is complementary to such serving techniques: unnecessary capacity that is never introduced does not need to be optimized, stored, transferred, or compressed later.

The problem is also scientifically important. Improvements reported by LoRA--MoE methods may originate from conditional specialization, but they may also result from a larger stored parameter count. Without a shared budget, these explanations cannot be separated. Moreover, different modules have unequal costs: one rank in a feed-forward projection can require far more parameters than one rank in a narrower attention projection. A nominally identical rank is therefore not an identical resource allocation. Establishing budget-matched evaluation is necessary for credible progress in PEFT.

Most importantly, adaptation demand is task-dependent but latent. The update that a module needs is unavailable until training is complete, yet completing training before choosing the architecture defeats the purpose of efficient allocation. A useful method must infer demand from limited evidence. This is difficult because activation magnitude, variation across inputs, effective dimensionality, and gradient sensitivity describe different properties. High energy does not necessarily imply high-dimensional demand, and high dimensionality does not necessarily imply that several experts will specialize. The allocator must make coupled structural decisions from noisy proxies while avoiding router collapse, redundant experts, and starvation of ordinary LoRA modules.

Solving this problem would turn adapter design from a manually tuned architecture choice into a measurable resource-allocation problem. It would enable fairer comparisons, lower the cost of adapting one model to many tasks, and provide a principled basis for deciding when conditional computation is worth its additional complexity. The topic therefore connects a timely theoretical question---how downstream data shapes useful parameter subspaces---with an immediate systems need for scalable specialization of foundation models.

\subsection{Method}

...


\end{document}
